\documentclass[11pt,a4paper]{article}

% =============================================================================
% Paquetes
% =============================================================================
\usepackage[utf8]{inputenc}
\usepackage[T1]{fontenc}
\usepackage[spanish,es-minimal]{babel}
\usepackage{amsmath,amssymb,amsthm}
\usepackage{geometry}
\geometry{margin=2.6cm}
\usepackage{booktabs}
\usepackage{longtable}
\usepackage{array}
\usepackage{enumitem}
\usepackage{tikz}
\usepackage{float}
\usepackage{hyperref}
\hypersetup{
  colorlinks=true,
  linkcolor=blue!60!black,
  urlcolor=blue!60!black,
  citecolor=blue!60!black,
  pdftitle={Las extensiones de HydroCouple al motor SWMM de la US EPA}
}
\usepackage{microtype}
\usepackage{xcolor}
\usepackage{siunitx}

% =============================================================================
% Entornos y macros
% =============================================================================
\newtheorem{definition}{Definición}[section]
\newtheorem{remark}{Observación}[section]

\newcommand{\st}{\mathrm{s}}
\newcommand{\ft}{\mathrm{ft}}
\newcommand{\cfs}{\mathrm{ft^3/s}}
\newcommand{\sfft}{\mathrm{ft^2}}
\newcommand{\R}{\mathbb{R}}
\renewcommand{\arraystretch}{1.25}

% Abreviatura para rutas del motor
\newcommand{\src}[1]{\texttt{src/engine/#1}}

\title{\textbf{Las Modificaciones de HydroCouple al Motor SWMM}\\[0.6em]
\large de la US EPA: el motor OpenSWMM y sus extensiones}
\author{Documento de Referencia Técnica}
\date{Agosto 2026}

\begin{document}

\maketitle

\begin{abstract}
Este documento describe, a nivel de ecuaciones y de diseño, el trabajo que el
grupo \textbf{HydroCouple} ha realizado sobre el motor de simulación
hidráulica \textbf{SWMM} de la \textbf{U.S.\ Environmental Protection Agency}
(US EPA). El motor \textbf{OpenSWMM} no es una modificación puntual del código
C del SWMM~5.x oficial: es una reescritura moderna en C$++$ que persigue dos
objetivos simultáneos: (i) \emph{paridad bit a bit} con el motor EPA para las
funcionalidades estándar, y (ii) una capa de \emph{extensiones numéricas y de
plataforma} que el motor oficial no posee. Este documento se enfoca en esa
segunda capa: la formulación semi-implícita de continuidad de nodo, la
aceleración de Anderson del ciclo de Picard, la ranura de Preissmann dinámica,
las uniones virtuales, el solver 1D explícito de volúmenes finitos, el módulo
2D con acoplamiento 1D--2D, la recuperación física del RDII, el paralelismo
OpenMP bit-exacto y los cambios de comportamiento por defecto.
\end{abstract}

\tableofcontents
\newpage

% =============================================================================
% 1. INTRODUCCIÓN
% =============================================================================
\section{Introducción}

\subsection{Contexto: OpenSWMM y su relación con SWMM de la US EPA}

El SWMM (\emph{Storm Water Management Model}) es el modelo de gestión de aguas
lluvias y aguas servidas de la US EPA. Su motor hidráulico 1D (SWMM~5.x)
resuelve las ecuaciones de de~Saint-Venant sobre una red de nodos y enlaces con
tres formulaciones de tránsito (escurrimiento permanente, onda cinemática y
onda dinámica), y ha evolucionado lentamente desde la década de 1970.

El proyecto \textbf{OpenSWMM} (motor \texttt{openswmm.engine}), liderado por
HydroCouple, reimplementa ese motor desde cero en C$++$ moderno con dos
principios de diseño:

\begin{enumerate}[nosep]
  \item \textbf{Paridad bit a bit} con el motor EPA para las funcionalidades
        estándar: cada término aritmético de los solvers clásicos reproduce
        exactamente el orden de operaciones y los valores del código legado
        (por ejemplo, el exponente $1.33333$ de la fricción de Manning, los
        factores de conversión truncados, la agrupación de operaciones en coma
        flotante). Esto se valida con una batería de pruebas de regresión
        contra el motor original y con \emph{soluciones manufacturadas}.
  \item \textbf{Una capa de extensiones}: algoritmos numéricos nuevos y
        opciones que el SWMM oficial no posee, orientados a mejorar la
        estabilidad, la convergencia y la fidelidad física de los transitorios,
        junto con arquitecturas de cómputo modernas (paralelismo, GPU,
        WebAssembly).
\end{enumerate}

Este documento se ocupa exclusivamente de la segunda capa. Las secciones
siguientes describen cada extensión con sus ecuaciones, su motivación y su
sintaxis de entrada. Todas las referencias de código apuntan al directorio
\texttt{src/engine/} del motor (\texttt{hydraulics/},
\texttt{hydraulics/fv/}, \texttt{2d/} y \texttt{hydrology/}).

\subsection{La reescritura arquitectónica como base}

Aunque no es una ``extensión'' funcional, la reescritura es el habilitador de
todo lo demás. Las decisiones arquitectónicas clave son:

\begin{itemize}[nosep]
  \item \textbf{Almacenamiento orientado a columnas (SoA)}: los datos de nodos
        y enlaces se almacenan como arreglos paralelos (\texttt{NodeData.hpp},
        \texttt{LinkData.hpp}), lo que permite accesos contiguos a caché y
        vectorización SIMD en los bucles internos.
  \item \textbf{Geometría por lotes}: la clase \texttt{XSectGroups}
        (\texttt{XSectBatch.cpp}) agrupa los conductos por forma geométrica y
        calcula áreas, radios hidráulicos y espejos de agua de todos los
        conductos en pasadas vectorizables, sin despacho por forma dentro del
        bucle interno.
  \item \textbf{Núcleos de momentum por categoría}: cada conducto se clasifica
        una vez por iteración en una categoría (seco, Manning a superficie
        libre, cerrado a sección llena, impulsión por Hazen--Williams o por
        Darcy--Weisbach) y se despacha a un núcleo sin ramas.
  \item \textbf{Paralelismo OpenMP bit-exacto}: el ciclo de Picard completo
        corre dentro de una única región \texttt{omp parallel} (equipo
        persistente), con un \texttt{THREADS n} en \texttt{[OPTIONS]}. La
        exactitud bit a bit a cualquier número de hilos se garantiza por
        diseño: cada bucle paralelo escribe elementos disjuntos (un solo
        productor por elemento), y el único combine entre hilos es un contador
        entero de nodos no convergidos, libre de orden. El barrido de
        caudales hacia los nodos se hace mediante una adyacencia CSR
        nodo$\to$conducto que reproduce exactamente el orden de acumulación en
        coma flotante del barrido serial (\texttt{buildConduitNodeCSR}).
\end{itemize}

% =============================================================================
% 2. CONTINUIDAD DE NODO SEMI-IMPLÍCITA
% =============================================================================
\section{Continuidad de nodo semi-implícita (\texttt{NODE\_CONTINUITY})}

\subsection{Motivación}

En el SWMM clásico, la actualización del tirante en un nodo usa dos ramas
separadas:
\begin{itemize}[nosep]
  \item \emph{Superficie libre}: $\Delta y = \Delta V / A_S$, con $A_S$ el área
        superficial del ensamble de nodo;
  \item \emph{En carga (EXTRAN)}: $\Delta y = \alpha\,\Delta Q / \mathrm{denom}$,
        donde el denominador incorpora el gradiente $\sum dQ/dH$.
\end{itemize}

La conmutación entre ambas ramas ocurre exactamente en la cota de la clave del
conducto más alto conectado al nodo. El operador de actualización de carga es,
por tanto, \emph{discontinuo} en la clave: un nodo que entra o sale de la
condición en carga cambia de fórmula de un paso al siguiente. Esta
discontinuidad es una fuente de oscilaciones y de iteraciones adicionales en
redes que transitan repetidamente por la clave (muy común en redes unitarias).

\subsection{La formulación unificada}

La extensión \texttt{NODE\_CONTINUITY SEMI\_IMPLICIT} reemplaza las dos ramas
por una única expresión suave. La idea es reconocer que los caudales que
entran a la actualización de la ecuación de continuidad dependen de la propia
carga que se está resolviendo. Linealizando el caudal neto alrededor de la
estimación actual con los gradientes de caudal:
\begin{equation}
  \sum Q^{t+\Delta t} \approx \sum Q + \sum\frac{dQ}{dH}\,\Delta H,
\end{equation}
y sustituyendo en la actualización trapezoidal de la carga,
\begin{equation}
  \Delta V = \tfrac{\Delta t}{2}\left(\sum Q^{t} + \sum Q^{t+\Delta t}\right),
\end{equation}
se obtiene la expresión unificada que el motor implementa en
\texttt{setNodeDepth} (\texttt{DynamicWave.cpp}):
\begin{equation}
  \Delta y = \frac{\Delta V}{A_S + \tfrac{\Delta t}{2}\sum \frac{dQ}{dH}},
  \qquad
  y^{t+\Delta t} = y^{t} + \Delta y,
\end{equation}
con el denominador pisado por el área mínima $A_{\min}$.

\subsection{Interpretación y convención de signo}

El motor acumula $\sum dQ/dH$ como una \emph{magnitud positiva}, con
$dQ_{\mathrm{net}}/dH = -\sum dQ/dH$ (una subida de carga aumenta la salida
neta). Bajo esa convención, el término $\tfrac{\Delta t}{2}\sum dQ/dH$ del
denominador es positivo y \emph{amortigua} la actualización: una subida de
carga que incrementa la salida neta reduce el llenado neto. El comportamiento
es el esperado en ambos regímenes:

\begin{itemize}[nosep]
  \item cuando el área superficial domina el denominador, la expresión colapsa
        a la fórmula clásica de superficie libre $\Delta y = \Delta V/A_S$;
  \item cuando el nodo se acerca y atraviesa la clave, el término de gradiente
        de caudal toma el rol que la fórmula de surcharge EXTRAN juega en la
        formulación explícita, sin cambiar de rama.
\end{itemize}

Esta es una \emph{corrección} sobre la ambigüedad de signo del manual del
SWMM oficial: el motor fue corregido para usar el signo más
(\texttt{DynamicWave.cpp:3374}, con el comentario de derivación
$\Delta H = \Delta V/(A_S + 0.5\,\Delta t\,\mathrm{sumdqdh})$ para
$\mathrm{sumdqdh} > 0$).

\subsection{Ventajas}

La continuidad de nodo semi-implícita:
\begin{itemize}[nosep]
  \item elimina la rama en la clave, por lo que el operador es $C^1$-suave a
        través de la transición superficie libre $\leftrightarrow$ en carga;
  \item es el ajuste \emph{recomendado} para modelos con uniones virtuales
        (Sección~\ref{sec:vj}), porque el acoplamiento de momentum a través del
        nodo requiere suavidad;
  \item es la configuración natural para combinarla con la aceleración de
        Anderson (Sección~\ref{sec:anderson}): bajo \texttt{EXPLICIT} los
        nodos en carga quedan excluidos de la aceleración, mientras que bajo
        \texttt{SEMI\_IMPLICIT} siguen siendo elegibles.
\end{itemize}

% =============================================================================
% 3. ACELERACIÓN DE ANDERSON
% =============================================================================
\section{Aceleración de Anderson del ciclo de Picard (\texttt{ANDERSON\_ACCEL})}
\label{sec:anderson}

\subsection{Motivación}

El solver de onda dinámica resuelve las ecuaciones de nodo y de enlace por
\emph{iteración funcional} (Picard / aproximaciones sucesivas): cada pasada
aplica el mismo operador de actualización de carga $G$ a las últimas
estimaciones de carga hasta que ningún nodo cambia más que la tolerancia de
carga $\varepsilon_H$. La iteración de punto fijo converge \emph{linealmente},
y el factor de subrelajación $\omega = 0.5$ que amortigua cada actualización
estabiliza la iteración sin mejorar su tasa. En redes con muchos nodos
fuertemente acoplados, el solver puede consumir sus 8 intentos
(\texttt{MAX\_TRIALS}) en casi todos los pasos de tránsito, incluso en
condiciones suaves.

\subsection{El método}

La aceleración de Anderson de profundidad 2 (equivalente a la actualización de
Aitken) mezcla las dos salidas más recientes del operador. Sea
$H_k$ la estimación de carga al entrar a la iteración $k$, el residual es
\begin{equation}
  r_k = G(H_k) - H_k,
\end{equation}
y el coeficiente de mezcla es
\begin{equation}
  \alpha_k = \mathrm{clamp}\!\left(
    \frac{r_k\,(r_k - r_{k-1})}{(r_k - r_{k-1})^{2}},\ 0,\ 1\right),
  \qquad
  H_{k+1} = (1-\alpha_k)\,G(H_k) + \alpha_k\,G(H_{k-1}).
\end{equation}
El acotamiento de $\alpha_k$ a $[0,1]$ restringe la actualización a una
\emph{interpolación} entre dos salidas del operador ya calculadas y ya acotadas;
nunca se produce una carga extrapolada. La mezcla se aplica por nodo a partir
del segundo intento de cada paso de tránsito.

\subsection{Salvaguardas}

\begin{itemize}[nosep]
  \item \textbf{Compuerta por magnitud del residual}: la mezcla se aplica solo
        cuando $\lvert r_k \rvert \le 20\,\varepsilon_H$; lejos del régimen
        lineal la mezcla podría sobrepasar.
  \item \textbf{Cotas físicas}: un tirante mezclado negativo se descarta en
        favor de la iteración ordinaria de Picard.
  \item \textbf{Exclusión de nodos no suaves}: la aceleración se omite en los
        nodos donde el operador $G$ es conocido no suave (Tabla~\ref{tab:aa}):
        nodos en carga bajo \texttt{EXTRAN} con continuidad \texttt{EXPLICIT}
        (la rama cambia en la clave), ranura dinámica activa, conductos cerca
        del corte de la ranura estática, vertederos u orificios en la clave,
        extremos de bombas (estado discreto on/off) y el borde de
        encharcamiento.
  \item \textbf{Criterio de convergencia doble}: un nodo se cuenta como
        convergido solo cuando tanto el residual crudo $\lvert G(H_k) - H_k
        \rvert$ como el movimiento aceptado $\lvert H_{k+1} - H_k \rvert$ están
        dentro de la tolerancia. Sin esta doble condición, una mezcla que
        aterrice cerca de la iteración anterior podría declarar convergencia
        mientras el balance de caudales sigue sin satisfacerse.
\end{itemize}

\begin{table}[H]
\centering
\small
\begin{tabular}{>{\raggedright}p{3.4cm}>{\raggedright}p{6.0cm}>{\raggedright\arraybackslash}p{5.4cm}}
\toprule
\textbf{Condición} & \textbf{Cuándo} & \textbf{Razón} \\
\midrule
Nodo en carga & \texttt{SURCHARGE\_METHOD EXTRAN} con
  \texttt{NODE\_CONTINUITY EXPLICIT} & la rama cambia en la clave \\
Ranura dinámica activa & \texttt{SURCHARGE\_METHOD DYNAMIC\_SLOT} & la geometría
  se reescribe por iteración \\
Cerca del corte estático & \texttt{SURCHARGE\_METHOD SLOT}, $0.98 \le
  \bar y/y_{\mathrm{full}} \le 1.02$ & el ancho de ranura entra abruptamente \\
Vertedero/orificio en la clave & LGH aguas arriba $\ge$ clave & la ecuación de
  caudal cambia de forma \\
Bombas & ambos nodos extremos & el estado on/off es discreto \\
Borde de encharcamiento & nodo encharcado en $d_{\mathrm{full}}$ & piso C$^0$ en
  el límite de encharcamiento \\
\bottomrule
\end{tabular}
\caption{Nodos excluidos de la aceleración de Anderson.}
\label{tab:aa}
\end{table}

El efecto medido: reducción de los conteos de intentos en aproximadamente 25 a
50\,\% por paso de tránsito en redes que de otro modo iteran hasta el límite.

% =============================================================================
% 4. RANURA DE PREISSMANN DINÁMICA
% =============================================================================
\section{Ranura de Preissmann dinámica}
\label{sec:dps}

\subsection{Motivación}

La opción \texttt{SURCHARGE\_METHOD DYNAMIC\_SLOT} agrega un tercer método de
escurrimiento en carga. El SWMM oficial ofrece dos formas de representar el
escurrimiento presurizado
en conductos cerrados: el algoritmo EXTRAN (perturbación tipo Hardy--Cross con
$dQ/dH$) y la ranura de Preissmann estática. Ambos tienen limitaciones
conocidas en los transitorios de llenado/vaciado: la ranura estática fija su
ancho a partir de la profundidad instantánea, lo que en el frente de mezcla
superficie libre/presión produce ``apretón de ranura'' (\emph{slot squeezing}),
una amplificación de energía artificial. La extensión \texttt{DYNAMIC\_SLOT}
implementa la formulación de ranura de Preissmann \emph{dinámica} de Sharior,
Hodges \& Vasconcelos (2023), en la que el área de la ranura evoluciona en el
tiempo como un elemento de \emph{almacenamiento transitorio}.

\subsection{Formulación}

La ranura dinámica tiene un ancho superior que depende del número de
Preissmann $P$ y de una celeridad de onda de presión objetivo $c_{pT}$
(el usuario la fija con \texttt{DPS\_CELERITY}, por defecto 25 m/s):
\begin{equation}
  T_s = \frac{g\,A_{\mathrm{full}}}{c_{pT}^{2}}\,P^{2},
\end{equation}
donde $A_{\mathrm{full}}$ es el área de sección llena. El número de Preissmann
parte, al inicio de la presurización, en
\begin{equation}
  P_{0} = \max\!\left(\frac{c_{pT}}{\alpha\,c_g},\ 1\right),
  \qquad c_g = \sqrt{g\,A_{\mathrm{full}}/W_{\max}},
\end{equation}
con $\alpha$ el parámetro de choque (\texttt{DPS\_ALPHA}, por defecto 3) y
$c_g$ la celeridad de la onda de gravedad a sección llena. El área de ranura se
acumula de forma \emph{dependiente de la trayectoria}:
\begin{equation}
  A_s \leftarrow \max\!\left(A_s + T_s\,\Delta h_s,\ 0\right),
  \qquad h_s = \max(\bar y - y_{\mathrm{full}},\ 0),
\end{equation}
es decir, cada incremento de almacenamiento se crea con el ancho de ranura
vigente en el momento en que se acumula, y las contribuciones previas nunca se
reescriben cuando $P$ decae. Esta propiedad es la que evita la amplificación de
energía del ``apretón de ranura''. Si la carga cae bajo la clave con área
residual, la carga de sobrepresión se mantiene en cero y el área restante
drena a través de incrementos negativos sucesivos (histéresis de
despresurización).

Después de la presurización, $P$ decae exponencialmente hacia 1 con la escala
de tiempo \texttt{DPS\_DECAY\_TIME} ($r = 0.5$ s por defecto):
\begin{equation}
  \hat P(t) = 1 + (\hat P_0 - 1)\exp\!\left(-\frac{10\,(t - t_s)}{r}\right),
\end{equation}
con $t_s$ el instante en que el conducto entró en carga, y se suaviza
espacialmente una vez por paso promediando $\hat P$ sobre los conductos
incidentes a cada nodo y tomando como $P$ de trabajo el promedio de los dos
extremos del conducto.

\subsection{Consecuencias para el solver}

Mientras una ranura está activa:
\begin{itemize}[nosep]
  \item el área hidráulica es $A = A_{\mathrm{full}} + A_s$ y el ancho superior
        es $T_s$;
  \item el área superficial aportada a un nodo extremo en carga es
        $T_s\,L/4$;
  \item el radio hidráulico se mantiene en su valor de sección llena (la
        ranura aporta almacenamiento pero no fricción);
  \item las cargas de los nodos se actualizan \emph{siempre} con la fórmula
        ordinaria de superficie libre; la rama de surcharge EXTRAN nunca se
        invoca, y la carga piezométrica sobre la clave emerge naturalmente como
        $z_{\mathrm{inv}} + y_{\mathrm{full}} + h_s$.
\end{itemize}
El paso de tiempo variable usa la celeridad de presión efectiva
$c_p = c_{pT}/P$ en la condición de Courant del conducto en carga:
\begin{equation}
  \Delta t \le \frac{L}{\lvert \bar U \rvert + c_{pT}/P}.
\end{equation}

% =============================================================================
% 5. UNIONES VIRTUALES
% =============================================================================
\section{Uniones virtuales (\texttt{[VIRTUAL\_JUNCTIONS]})}
\label{sec:vj}

\subsection{Motivación}

Un cambio de pendiente en un conducto debe dividirse en una cámara de unión.
Pero una cámara de unión ordinaria introduce dos artefactos numéricos:
\begin{enumerate}[nosep]
  \item \textbf{Almacenamiento artificial}: su área superficial se pisa con el
        mínimo $A_{\min} = 12.566\ \sfft$, lo que agrega un volumen de
        estancamiento que ``emborrona'' los transitorios;
  \item \textbf{Ruptura de momentum}: el nodo actúa como un pequeño volumen de
        estancamiento que interrumpe la transmisión de momentum entre los dos
        conductos.
\end{enumerate}

Una \textbf{unión virtual} (\texttt{[VIRTUAL\_JUNCTIONS]}, sección con nombre y
cota de batea; todo lo demás se deriva) elimina ambos artefactos para dos
conductos colineales de sección idéntica que se encuentran en un quiebre de
pendiente. Se declara un nodo sellado de almacenamiento idénticamente nulo. Las
reglas de elegibilidad se verifican en la lectura: exactamente dos enlaces,
ambos conductos; secciones transversales idénticas (forma, dimensiones,
referencia de curva, número de barriles; la rugosidad de Manning puede
diferir); ambos desfases en el nodo nulos con batea continua; sin aportes
laterales de ningún tipo; y método de tránsito de onda dinámica. Cualquier
violación es un error de entrada.

\subsection{Tratamiento de continuidad}

Una unión virtual es un nodo sellado de almacenamiento cero. Su carga se
actualiza con la fórmula de superficie libre usando el área superficial natural
de medio enlace aportada por sus dos conductos, \emph{sin} el piso de área
mínima (ese piso es precisamente el almacenamiento artificial que la
funcionalidad elimina). Cuando el área natural se anula (par seco, o par en
carga con ancho de ranura pequeño), la actualización cae a una forma pura de
balance de caudales de la actualización en carga con $\alpha = 1$ y sin piso.
En la convergencia, $\sum Q = 0$ en el nodo. El nodo está sellado: su carga
puede subir sobre la clave sin límite (como una cámara con tapa empernada),
nunca se anega ni se encharca, y su volumen y rebalse son idénticamente cero.

\subsection{Tratamiento de momentum}

Para un par de paso (un conducto entrando, otro saliendo), el conducto aguas
abajo toma como estado aguas arriba los valores medios del conducto aguas
arriba en su ponderación por Froude (mecanismo de \emph{upwinding} a través
del nodo), transportando el momentum advectado a través del nodo en lugar de
reiniciarlo. Con \texttt{VIRTUAL\_JUNCTION\_MOMENTUM FULL} (por defecto
\texttt{BASIC}) se agrega además una corrección convectiva al término de
inercia de ambos conductos:
\begin{equation}
  \Delta Q_j = \Delta t\,\sigma_j\,
    \frac{\left(\bar U^{2}\bar A\right)_{\mathrm{ab}} -
          \left(\bar U^{2}\bar A\right)_{\mathrm{ar}}}{\Lambda},
  \qquad \Lambda = \frac{L_{\mathrm{ar}} + L_{\mathrm{ab}}}{2},
\end{equation}
con $\sigma_j$ un factor de amortiguación de la forma de Froude; la corrección
se anula cuando los dos conductos no transportan caudal en el mismo sentido a
través del nodo. Los pares en silla (ambos conductos hacia el nodo) o en cresta
(ambos desde el nodo) reciben el tratamiento de continuidad de almacenamiento
cero pero no el acoplamiento direccional de momentum. La pareja siempre se
resuelve junta (no se congela por el bypass de convergencia) y se agrega una
verificación de Courant a nivel de par.

La guía de uso: las uniones virtuales están pensadas para quiebres de
pendiente de pequeña deflexión; los cambios de alineación horizontal deben
seguir siendo cámaras de unión ordinarias con coeficientes de pérdida.

% =============================================================================
% 6. SOLVER 1D EXPLÍCITO DE VOLÚMENES FINITOS
% =============================================================================
\section{Tránsito 1D por volúmenes finitos explícitos (\texttt{FLOW\_ROUTING FV})}
\label{sec:fv}

\subsection{Motivación}

El SWMM oficial solo ofrece esquemas implícitos con linealización (onda
dinámica) u ondas cinemática/permanente. El solver de onda dinámica, aunque es
el estándar de la industria, tiene dificultades bien documentadas con el flujo
transcrítico, la propagación de frentes bruscos y la transición seco/húmedo. La
extensión \texttt{FLOW\_ROUTING FV} agrega un solver 1D de \emph{volúmenes
finitos explícito y conservativo} de tipo Godunov, que resuelve el flujo
transcrítico de forma nativa y maneja seco/húmedo sin las ramas ad hoc de la
onda dinámica.

\subsection{Arquitectura del solver}

El solver vive en \src{hydraulics/fv/}. Cada conducto se discretiza en celdas
(por defecto al menos \texttt{FV\_MIN\_CELLS} = 4 por conducto, o con un blanco
de longitud \texttt{FV\_CELL\_LENGTH}); el estado conservado de cada celda es
el área hidráulica $A$ y el caudal $Q$. El paso temporal global obedece la
condición de Courant, y el solver sub-pasea internamente a su propio límite
CFL, de modo que el paso de tránsito (de reporting/forzamiento) no necesita
encogerse por estabilidad.

Las opciones principales (claves \texttt{FV\_*} en \texttt{[OPTIONS]}) son:

\begin{table}[H]
\centering
\small
\begin{tabular}{>{\raggedright}p{4.6cm}>{\raggedright}p{2.6cm}>{\raggedright\arraybackslash}p{7.2cm}}
\toprule
\textbf{Opción} & \textbf{Defecto} & \textbf{Descripción} \\
\midrule
\texttt{FV\_CFL} & 0.5 & número de Courant \\
\texttt{FV\_RIEMANN} & HLLC & solver de Riemann (HLL / HLLC) \\
\texttt{FV\_ORDER} & 1 & orden espacial (1 = Godunov, 2 = MUSCL--Hancock) \\
\texttt{FV\_LIMITER} & MINMOD & limitador de pendiente (MINMOD/VANLEER/SUPERBEE) \\
\texttt{FV\_CELL\_LENGTH} & 0 & blanco de $\Delta x$ (0 = solo piso de celdas) \\
\texttt{FV\_MIN\_CELLS} & 4 & piso de celdas por conducto \\
\texttt{FV\_SLOT\_CELERITY} & 100 ft/s & celeridad de presión (ancho de ranura) \\
\texttt{FV\_NODE\_COUPLING} & SEMI\_IMPLICIT & acoplamiento nodo--celda \\
\texttt{FV\_STRUCTURE\_COUPLING} & SUBSTEP & cuándo se re-evalúan las estructuras \\
\texttt{FV\_LTS} & true & paso de tiempo local \\
\texttt{FV\_LTS\_MAX\_TIERS} & 6 & máximo de niveles LTS \\
\texttt{FV\_DISPERSION} & 0 & coeficiente de dispersión longitudinal \\
\bottomrule
\end{tabular}
\caption{Opciones del solver 1D de volúmenes finitos.}
\label{tab:fv}
\end{table}

\subsection{Diseños clave}

\begin{itemize}[nosep]
  \item \textbf{Nodos algebraicos}: las cámaras de unión no son estados con
        área (como en la onda dinámica, donde el área de trabajo pertenece a
        los conductos); un nodo de grado 2 sin aporte lateral pasa los flujos
        directamente como una cara, y los demás nodos resuelven su carga desde
        el balance instantáneo de flujos por subpaso. Esto elimina el límite de
        paso de milisegundos que un área de cámara imponía al modelo completo.
  \item \textbf{Acoplamiento nodo--celda semi-implícito}: el flujo de cada cara
        de acoplamiento se linealiza en la carga del nodo, lo que saca al nodo
        del límite de estabilidad explícito (la cámara, no la tubería, era el
        factor limitante del subpaso). La conservación no se altera: el flujo
        corregido es el que ven tanto el nodo como la celda. Con
        \texttt{FV\_NODE\_PICARD\_SWEEPS} $> 1$ se re-evalúan área, ancho y
        flujos de Riemann en cada barrido.
  \item \textbf{Paso de tiempo local (LTS)}: las celdas rígidas (conductos
        cortos, presurizados) avanzan a su propio $\Delta t = 2^k\,\Delta t_0$
        mientras el resto avanza al paso global; cuando el reparto no separa
        nada, el solver cae al camino de paso global bit a bit.
  \item \textbf{Estructuras}: las ecuaciones de bombas/orificios/vertederos/
        descargas se re-evalúan por defecto en \emph{cada} subpaso explícito
        (\texttt{FV\_STRUCTURE\_COUPLING SUBSTEP}), lo que es físicamente
        exacto; la alternativa es una vez por paso de tránsito.
  \item \textbf{Backends}: el mismo núcleo escalar se compila para CPU,
        OpenMP y dispositivos (CUDA/HIP/SYCL mediante plugins Kokkos); el
        resultado debe ser bit a bit idéntico entre backends.
\end{itemize}

% =============================================================================
% 7. MÓDULO 2D Y ACOPLE 1D–2D
% =============================================================================
\section{El módulo 2D y el acoplamiento 1D--2D}
\label{sec:2d}

La extensión de mayor envergadura es el módulo de \emph{escurrimiento
superficial 2D}: una malla triangular no estructurada sobre la que se resuelven
las ecuaciones de aguas someras en su forma \emph{local-inercial}
(de~Almeida \& Bates, 2013):
\begin{equation}
  \frac{\partial h}{\partial t} + \nabla\cdot \mathbf{q} = i - e + s,
  \qquad
  \frac{\partial \mathbf{q}}{\partial t} + g\,h\,\nabla\eta +
  g\,n^{2}\,\frac{\lvert\mathbf{q}\rvert\,\mathbf{q}}{h^{7/3}} = 0,
\end{equation}
donde $h$ es el tirante, $\mathbf{q}$ el caudal específico (m$^2$/s),
$\eta$ la cota de superficie libre, $n$ la rugosidad de Manning, $i$ la lluvia
sobre la malla y $e$, $s$ los términos de evaporación y fuente. Se conserva el
término $\partial \mathbf{q}/\partial t$ (inercia local) pero se omite la
aceleración convectiva: preciso para flujo subcrítico de pendiente suave
($\mathrm{Fr} \lesssim 0.5$). El módulo se activa por la sola presencia de una
malla (\texttt{[2D\_VERTICES]} / \texttt{[2D\_TRIANGLES]} o archivo de malla),
y su solver está portado a GPU/WebGPU con exactitud bit a bit respecto del
backend serial.

\subsection{La ley de orificio de acoplamiento}

El intercambio entre un nodo 1D y una celda 2D es una ley de descarga tipo
orificio gobernada por la diferencia de carga $\Delta h = h_{2D} - h_{1D}$
(\texttt{NodeCoupling.cpp}):
\begin{equation}
  Q = C_d\,A_{\mathrm{eff}}\,\mathrm{sign}(\Delta h)\,\sqrt{2g}\,
      \varphi(\lvert\Delta h\rvert),
\end{equation}
positiva en el sentido 2D$\to$1D (drenaje). La raíz cuadrada se regulariza
con la función $C^1$
\begin{equation}
  \varphi(x) = \begin{cases}
    \dfrac{3}{2\sqrt{\varepsilon}}\,x - \dfrac{1}{2\varepsilon^{3/2}}\,x^{2}
      & x < \varepsilon,\\[1em]
    \sqrt{x} & x \ge \varepsilon,
  \end{cases}
  \qquad \varepsilon = 0.02\ \mathrm{m},
\end{equation}
que coincide con $\sqrt{x}$ en valor y pendiente en $\varepsilon$ y tiene
pendiente finita en el origen. Esto elimina la singularidad
$dQ/d\Delta h \to \infty$ en $\Delta h \to 0$, exactamente el régimen de
equilibrio llenado/vaciado donde el orificio desnudo era una fuente de
rigidez y oscilaciones.

Además se aplican tres regularizaciones:
\begin{itemize}[nosep]
  \item \textbf{Compuerta de conducto con tapa}: sobre la clave, una
        \emph{smoothstep} de Hermite en una banda de 5 cm multiplica $Q$; bajo
        la clave el intercambio es exactamente cero (el nodo solo intercambia
        cuando sobrepasa la clave). El área efectiva crece linealmente de $A$
        en la clave a $2A$ a 5 cm sobre ella.
  \item \textbf{Rampa seco/húmedo del lado fuente}: $Q$ se multiplica por una
        smoothstep sobre el tirante vivo del lado fuente relativo al
        \texttt{DRY\_DEPTH}, de modo que un drenaje se autolimita a cero cuando
        la celda se vacía y un rebalse se autolimita cuando el nodo se vacía.
  \item \textbf{Sensibilidad a la carga (Jacobiano)}: se dispersa la derivada
        $-dQ/dh_{1D}$ en el denominador de continuidad del nodo de la onda
        dinámica ($\sum dQ/dH$), convirtiendo la fuente explícita en un
        amortiguador que estabiliza la iteración de drenaje/rebalse.
\end{itemize}

\subsection{Presupuestos de volumen y cadencia}

El intercambio se evalúa \emph{en vivo} dentro del marchante 2D en la cadencia
de nivel 0 del paso de tiempo local, contra cargas 2D actuales y cargas 1D
congeladas al inicio del lote, acumulando $\int Q\,dt$ exactamente por punto.
Dos topes duros garantizan que un drenaje no pueda tomar más de una fracción
$\beta$ de la celda por subpaso y que la misma agua no pueda rebalsar dos veces
dentro de un paso de tránsito. Los volúmenes de intercambio de las cámaras se
entregan al 1D a través de una \emph{cola de entrega} que drena a tasa uniforme
sobre el lapso del lote (no como un pulso $\mathrm{volumen}/\Delta t$ que
inundaba cámaras pequeñas y provocaba oscilaciones).

\subsection{Acoplamiento por área de encharcamiento}

Para que el HGL de un nodo acoplado pueda subir sobre la clave y rastrear la
superficie 2D (y para que el rebalse/ingreso dispare), el router 2D
\emph{sobrescribe el área de encharcamiento} del nodo con la huella de las
celdas 2D circundantes (el área dual mediana de las celdas del stencil), y
marca al nodo como capaz de encharcarse \emph{independientemente} de la opción
global \texttt{ALLOW\_PONDING}. Esto es el \texttt{nodeCanPond} de la onda
dinámica extendido a nodos acoplados: el área de encharcamiento no es un dato
del usuario sino el área de la celda 2D asignada.

% =============================================================================
% 8. RECUPERACIÓN FÍSICA DEL RDII
% =============================================================================
\section{Recuperación física de la abstracción del RDII (\texttt{[RDII\_DECAY]})}

\subsection{Motivación}

En el SWMM oficial, la infiltración/entrada por lluvia (RDII, \emph{rain-derived
inflow and infiltration}) se modela con el método RTK: cada nodo responde con
hasta tres hidrogramas unitarios triangulares (corto, medio y largo) definidos
por $R$, $T$ y $K$, convolucionados con la precipitación. La abstracción
inicial (capacidad de almacenamiento antes de que se genere RDII) se recupera
entre eventos a \emph{tasa constante} dada por una tabla mensual
(\texttt{IA\_Recov}): la variación estacional debe pre-cocinarse en los datos de
entrada del usuario.

\subsection{La extensión: abstracción de decaimiento exponencial}

La extensión \texttt{[RDII\_DECAY]} reemplaza esa recuperación constante por un
modelo físico de \emph{relajación exponencial de primer orden} cuya tasa depende
de la \emph{temperatura del aire}. Los parámetros por par (grupo de
hidrogramas, respuesta) son $k_{\mathrm{dep}}$, $k_0$, $k_T$, $T_{\mathrm{ref}}$,
$\theta_{\mathrm{rec}}$, $T_{\mathrm{cong}}$ y, opcionalmente, un particionado
lluvia/nieve.

\textbf{Agotamiento durante la lluvia} ($k_{\mathrm{dep}}$): la capacidad
disponible se agota exponencialmente con la lámina de lluvia $\Delta P$,
\begin{equation}
  IA_{\mathrm{avail}}^{+} = IA_{\mathrm{avail}}\,
    e^{-k_{\mathrm{dep}}\,\Delta P},
  \qquad
  P_{\mathrm{net}} = \max\!\left(0,\ \Delta P -
    \left(IA_{\mathrm{avail}} - IA_{\mathrm{avail}}^{+}\right)\right),
\end{equation}
con contabilidad másica consistente (el almacenamiento drena exactamente lo
que abstrae). $k_{\mathrm{dep}} = 0$ desactiva la abstracción.

\textbf{Recuperación entre eventos} ($k_0$, $k_T$, $T_{\mathrm{ref}}$): la
capacidad disponible relaja hacia el máximo,
\begin{equation}
  \frac{d\,IA_{\mathrm{avail}}}{dt} = k_{\mathrm{rec}}(T)\,
    \left(IA_{\max} - IA_{\mathrm{avail}}\right),
  \qquad
  IA_{\mathrm{avail}}^{+} = IA_{\max} - \left(IA_{\max} - IA_{\mathrm{avail}}\right)
    e^{-k_{\mathrm{rec}}(T)\,\Delta t},
\end{equation}
donde la tasa de recuperación depende de la temperatura con una base
independiente de la temperatura (drenaje gravitacional / redistribución
capilar) más un término térmico:
\begin{equation}
  k_{\mathrm{rec}}(T) = \max\!\left(0,\ k_0 + k_T\,
    e^{\theta_{\mathrm{rec}}\,(T - T_{\mathrm{ref}})}\right),
\end{equation}
y se anula bajo la temperatura de congelamiento $T_{\mathrm{cong}}$ (suelo
congelado), reproduciendo la elevación de RDII de inicios de primavera. La
temperatura proviene de la fuente \texttt{[TEMPERATURE]} del proyecto (serie de
tiempo o archivo); si no se configura, $T$ se fija en $T_{\mathrm{ref}}$ y se
emite una advertencia.

\textbf{Particionado lluvia/nieve (opcional)}: con la cláusula \texttt{SNOW},
$T \le T_{\mathrm{nieve}}$ acumula la precipitación como equivalente en nieve
(sin aporte líquido), y con nieve presente y $T > T_{\mathrm{nieve}}$ se agrega
un derretimiento de grado-día $m = \min(SWE,\ DDF\,(T - T_{\mathrm{nieve}})\,
\Delta t)$ a la lluvia (lluvia sobre nieve).

La recuperación de primer orden es más rápida cuando el déficit es grande y se
ralentiza cerca de la saturación, a diferencia de la tasa constante lineal del
SWMM oficial. Los pares sin fila en \texttt{[RDII\_DECAY]} usan el modelo
lineal legado; el estado de ejecución (el tirante $ia_{\mathrm{used}}$) es
intercambiable entre formulaciones, por lo que los archivos de arranque en
caliente son compatibles.

% =============================================================================
% 9. CAMBIOS DE COMPORTAMIENTO Y PLATAFORMA
% =============================================================================
\section{Cambios de comportamiento por defecto y de plataforma}

\subsection{Cambios de comportamiento}

Además de las extensiones funcionales, el motor introduce cambios de
comportamiento que conviene conocer:

\begin{itemize}[nosep]
  \item \textbf{Paso de tiempo variable por defecto}: \texttt{VARIABLE\_STEP}
        por defecto es \textbf{0.75} (paso adaptativo por Courant), mientras
        que el manual del SWMM 5.x documenta 0 (paso fijo). Para reproducir el
        comportamiento del SWMM oficial debe fijarse \texttt{VARIABLE\_STEP 0}.
  \item \textbf{Manejo de unidades de las tolerancias}: la \texttt{HEAD\_TOLERANCE}
        y el \texttt{MIN\_SURFAREA} se convierten desde las unidades del
        proyecto al sistema interno (pies) al inicializar, corrigiendo errores
        de factor 3.3--10.8 en modelos SI del motor anterior.
  \item \textbf{Pérdidas de conducto por iteración}: bajo onda dinámica, la
        evaporación y filtración de conductos se recalculan en \emph{cada}
        iteración de Picard (con la compuerta de clase de escurrimiento),
        igual que el motor legado, en lugar de una vez por paso.
  \item \textbf{Recuento de convergencia}: un paso que converge en su último
        intento permitido se cuenta como convergido (solo falla si no converge
        tras \texttt{MAX\_TRIALS}).
  \item \textbf{Salto de estado estacionario}: \texttt{SKIP\_STEADY\_STATE}
        omite el tránsito cuando no hubo acciones de control, el error de
        continuidad está bajo la tolerancia y los aportes no cambiaron.
\end{itemize}

\subsection{Extensiones de plataforma y formato}

\begin{itemize}[nosep]
  \item \textbf{CRS}: la opción \texttt{CRS} en \texttt{[OPTIONS]} acepta un
        EPSG o una cadena PROJ, transportando el sistema de referencia espacial
        del modelo (el SWMM oficial no la posee).
  \item \textbf{Claves de extensión}: \texttt{ext\_options} almacena
        cualquier clave desconocida de \texttt{[OPTIONS]} para plugins.
  \item \textbf{IGNORE\_2D}: permite desactivar el solver 2D conservando la
        malla en el archivo (para ejecutar el modelo solo 1D).
  \item \textbf{Geopackage}: entrada/salida nativa en GeoPackage
        (incluida la tabla \texttt{rdii\_decay} de los parámetros de
        recuperación), además del formato \texttt{.inp}.
  \item \textbf{API C}: una API en C (con envoltorios) expone todas las
        extensiones (por ejemplo \texttt{swmm\_rdii\_decay\_add}).
  \item \textbf{WebAssembly}: el motor se compila a WASM para ejecutarse en el
        navegador (como usa LocalSWMM), con el solver 2D en GPU/WebGPU.
  \item \textbf{Pruebas de soluciones manufacturadas}: el repositorio incluye
        una batería de benchmarks con soluciones analíticas (rotura de presa
        de Ritter, lago en reposo, ondas de macdonald, curvas de referencia de
        impulsiones, entre otras) que anclan la fidelidad de los solvers.
\end{itemize}

% =============================================================================
% 10. RESUMEN DE OPCIONES NUEVAS
% =============================================================================
\section{Resumen de opciones nuevas del motor}

La Tabla~\ref{tab:nuevas} resume las opciones y secciones que HydroCouple
agrega sobre el SWMM oficial.

\begin{longtable}{>{\raggedright}p{5.4cm}>{\raggedright}p{5.0cm}>{\raggedright\arraybackslash}p{3.8cm}}
\toprule
\textbf{Opción / sección} & \textbf{Qué agrega} & \textbf{Valor por defecto} \\
\midrule
\endhead
\texttt{NODE\_CONTINUITY} & continuidad de nodo unificada semi-implícita &
  EXPLICIT (legado) \\
\texttt{ANDERSON\_ACCEL} & aceleración de Anderson del ciclo de Picard & NO \\
\texttt{SURCHARGE\_METHOD} & tercer método de carga: DYNAMIC\_SLOT &
  EXTRAN (legado) \\
\texttt{DPS\_CELERITY} / \texttt{DPS\_ALPHA} / \texttt{DPS\_DECAY\_TIME} &
  parámetros de la ranura dinámica & 25 m/s, 3, 0.5 s \\
\texttt{[VIRTUAL\_JUNCTIONS]} & nodos sellados de almacenamiento cero con
  transmisión de momentum & --- \\
\texttt{VIRTUAL\_JUNCTION\_MOMENTUM} & corrección convectiva a través de la
  unión virtual & BASIC \\
\texttt{FLOW\_ROUTING FV} & solver 1D explícito de volúmenes finitos &
  DYNWAVE (legado) \\
\texttt{FV\_*} & opciones del solver FV (CFL, Riemann, orden, LTS, ...) &
  ver Tabla~\ref{tab:fv} \\
\texttt{[2D\_VERTICES]} \\[-2pt] \texttt{[2D\_TRIANGLES]} & malla 2D de aguas
  someras local-inercial & --- \\
\texttt{[2D\_OPTIONS]} & opciones del marchante 2D y del acople
  (\texttt{COUPLING\_CD}, \texttt{COUPLING\_SYNC}, ...) & COUPLING\_CD 0.65 \\
\texttt{IGNORE\_2D} & desactiva el solver 2D conservando la malla & NO \\
\texttt{[RDII\_DECAY]} & recuperación exponencial de la abstracción del RDII
  con temperatura ($k_0$, $k_T$, $T_{\mathrm{ref}}$) & --- \\
\texttt{THREADS} & hilos OpenMP bit-exactos del solver & 1 \\
\texttt{CRS} & sistema de referencia espacial (EPSG/PROJ) & vacío \\
\texttt{SKIP\_STEADY\_STATE} & omite el tránsito en régimen permanente & NO \\
\bottomrule
\caption{Opciones y secciones nuevas del motor OpenSWMM respecto del SWMM
oficial de la US EPA.}
\label{tab:nuevas}
\end{longtable}

% =============================================================================
% 11. REFERENCIAS
% =============================================================================
\section{Referencias}

\begin{enumerate}[nosep,label={[\arabic*]}]
  \item Rossman, L.~A. (2017). \emph{Storm Water Management Model Reference
        Manual Volume II -- Hydraulics}. U.S.\ EPA.
  \item Sharior, S., Hodges, B.~R., \& Vasconcelos, J.~G. (2023).
        Generalized, dynamic, and transient-storage form of the Preissmann
        slot. \emph{Journal of Hydraulic Engineering}, 149(11).
  \item de~Almeida, G.~A.~M., \& Bates, P.~D. (2013). Applicability of the
        local inertial approximation of the shallow water equations to flood
        modeling. \emph{Water Resources Research}, 49(8).
  \item Código fuente del motor OpenSWMM (HydroCouple):
        \begin{itemize}[nosep,itemsep=0pt]
          \item \src{hydraulics/DynamicWave.cpp} \\
                \src{hydraulics/HydStructures.cpp} \\
                \src{hydraulics/fv/ExplicitFvSolver.cpp}
          \item \src{engine/2d/} (marchante y acople) \\
                \src{engine/hydrology/RDII.cpp} \\
                \src{core/SimulationOptions.hpp}
        \end{itemize}
  \item Manuales de referencia del motor OpenSWMM (Vol.\ II -- Hydraulics,
        cap.\ 3, 7, 9; Vol.\ I -- Hydrology, cap.\ 7).
\end{enumerate}

\end{document}
